\section{Recovery is one mechanism}
\label{sec:ha:recovery_shape}

\Cref{sec:ha:disabled} asks the same four questions of every component. The
answers to the third of them --- how a restarted instance becomes current --- have
a common shape, which is set out here because the next section depends on it.

\textbf{Every component recovers the same way.} Whatever it holds and whoever
holds the authority for it, the operation decomposes into four steps:

\begin{enumerate}[nosep]
  \item establish a \emph{position}: a point up to which its state is current;
  \item present that position to whatever holds the authoritative record;
  \item receive everything after it and apply that silently, without repeating
        side effects that were performed the first time;
  \item begin acting.
\end{enumerate}

What differs between components is only where the position in step one comes
from.

\begin{center}
\begin{tabular}{lll}
\toprule
\textbf{Component} & \textbf{Position, on a restart} & \textbf{Position, on a promotion} \\
\midrule
Authentication service & how current its credentials are & no promotion \\
Gateway                & the sequencer, on binding       & no promotion \\
Sequencer              & its log, and gateways reporting & the replica it maintained \\
Matching engine        & a checkpoint on disk            & the replica it maintained \\
Publisher              & each subscriber presents its own & each subscriber presents its own \\
Recorder               & its own confirmed publish        & its own confirmed publish \\
\bottomrule
\end{tabular}
\end{center}

Three things follow from this.

\textbf{Promotion is not a second mechanism.} For the components that have one,
promotion is this same operation with the position learned from a replica rather
than from a file or from being told. \Cref{sec:ha:enabled} is therefore not
another recovery design; it is about how an instance learns it may act, and about
the failures that exist only because there are two of them.

\textbf{The mechanism is exercised constantly rather than rarely.} A recovery
path used only on promotion runs only when a scenario contrives one. The same
path, used on every restart, runs whenever anything is upgraded, reconfigured or
dies --- which is the difference between a mechanism that is known to work and
one that is believed to.

\textbf{Its requirements are shared, and are stated once.} Because the operation
is common, so are the properties it needs, and they are not repeated per
component: \reqref{R-0031}, that no component acts on partially recovered state;
\reqref{R-0032}, that an interrupted recovery can be repeated without applying
anything twice; and \reqref{R-0064}, that the generation stamped on what is
served never goes backwards.

\textbf{Retention is the obligation every authoritative source shares.} A
position is only useful if what follows it still exists, so each source ---
the sequencer's log, the publisher's log, a member's reports --- must retain
records at least as far back as the oldest position any party may still present
to it. That bound is what \reqref{R-0008}, \reqref{R-0022} and \reqref{R-0048}
state for each source in turn. It is a single policy stated three times, and in
each case the retained period must be a figure the venue can state rather than
whatever storage happens to hold.

\begin{requirement}{R-0101}{A component establishes that its catch-up was complete before it acts}
  A component shall establish that it received every record between the position
  it presented and the position it has reached, and shall not begin acting if it
  did not.
  \rationale{Catch-up records are applied without producing the execution reports
  that were produced the first time, so nothing about applying them draws
  attention. A catch-up that is missing a record therefore leaves a component
  whose state is wrong and which believes it is current --- and every check it
  performs afterwards is made against the wrong state. Silence is the right
  behaviour for records that arrive; it is not the right behaviour for records
  that do not.}
  \uncovered{no scenario withholds a record during a catch-up}
\end{requirement}

\textbf{Three components depart from this shape, and the differences matter.} A
gateway and an
authentication service are never promoted, so they have one operation instead of
two --- which is why neither needs an arbiter, and why running several of each is
cheap. And a restarted sequencer recovers \emph{less} than a promoted one,
because a session that was not bound at the time is one nobody will report; that
asymmetry is the reason for \reqref{R-0012}.
